%% toc: P15 158 158.3,158.4
% Integration supplements from Durrett, secs. 1.5-1.6: monotone limits with an integrable lower part, signed series, convergence in measure, Jensen's consequences, moment bounds, inclusion-exclusion and Bonferroni.
\chapter{Integration Supplements: Limits, Jensen and Moment Bounds}
\label{sec:p15-w02c-p15-005-integration-supplements}

Chapter~\ref{sec:p15-integration-expectation} constructs the integral
and proves its main limit theorems and inequalities. This chapter
proves the further statements that Durrett's Sections~1.5 and~1.6 make
in passing or set as exercises \cite[pp.~26--37]{DurrettPTEAuthor2019}.
They answer six practical questions. When may a limit pass through an
integral of functions that are not nonnegative? When may a signed series
be integrated term by term? How do almost-everywhere convergence and
convergence in measure compare on a finite measure space? Which way
does Jensen's inequality point, and when is it an equality? What do a
mean and a variance say about a tail probability, and what can they not
say? How are the probabilities of intersections of events turned into
bounds on the probability of their union?

\emph{Inputs.} A measure space \((\Omega,\mathcal F,\mu)\), with no
finiteness assumed unless a statement says so; a probability space
\((\Omega,\mathcal F,\mathbb P)\) for the expectation statements; real
or extended real measurable functions; random variables with the
moments named in each hypothesis; and finitely many events.
\emph{Conclusions.} Monotone convergence for increasing sequences whose
first term has an integrable negative part; term-by-term integration of
an absolutely integrable signed series; continuity of an indefinite
Lebesgue integral; the implication from almost-everywhere convergence to
convergence in measure on a finite measure space, with its limits and
its use in bounded convergence; Jensen's inequality on \(\mathbb R^n\), its equality case, its two
special cases and the arithmetic--geometric mean inequality; the sharp
and the asymptotic behaviour of Chebyshev's bound, the one-sided bound,
the impossibility of a lower bound from a mean and a variance, and a
lower bound for \(\mathbb P(Y>0)\); a sequence for which a moment bound
gives convergence of means although no integrable dominator exists;
the behaviour of \(y\,\mathbb E(1/X;X>y)\) at both ends; and
inclusion--exclusion with the alternating Bonferroni bounds.

\emph{Definitions and assumptions.} The integral, signed integral,
integrability, convergence in measure, expectation and variance are those
of \cref{def:p15-nonnegative-integral,def:p15-integrable-functions,def:p15-lp-convergence,def:p15-expectation-integral},
and convexity on the line is \cref{def:p15-convex-integral}. Durrett
writes \(E\), \(P\) and \(\operatorname{var}\) for expectation,
probability and variance; here they are \(\mathbb E\), \(\mathbb P\)
and \(\operatorname{Var}\), as in the chapter this one extends. Lebesgue
measure \(\lambda\) and the uniform probability space
\(((0,1),\mathcal F_U,\mathbb U)\) are those of
\cref{thm:p15-lebesgue-unit-interval}. Durrett states his integration
results for sigma-finite measures \cite[p.~18]{DurrettPTEAuthor2019};
no proof below uses sigma-finiteness.

\section{Monotone limits, signed series and indefinite integrals}
\label{sec:p15-w02c-p15-005-limits-series}

\Cref{thm:p15-monotone-integral} requires nonnegative functions, and
the example \(r_n=-\mathbf1_{[n,\infty)}\) of
\cref{def:p15-convergence-failures} shows that monotonicity alone does
not suffice. The first result shows exactly what replaces
nonnegativity: an integrable negative part of the first function.

\begin{hbtheorem}{thm:p15-w02c-p15-005-monotone-lower}{proved=annex:p15-w02c-p15-005-monotone-lower}{Monotone convergence with an integrable negative part}
\par\nobreak
Let \(g_n,g\) be extended real measurable functions with
\(g_n\leq g_{n+1}\) and \(g_n\to g\) a.e., and suppose
\(\int g_1^-\,d\mu<\infty\). Then every \(\int g_n\,d\mu\) and
\(\int g\,d\mu\) exists in \((-\infty,\infty]\), and
\[
 \int g_n\,d\mu\ \uparrow\ \int g\,d\mu .
\]
On a probability space this reads: if \(\mathbb EX_1^-<\infty\) and
\(X_n\uparrow X\), then \(\mathbb EX_n\uparrow\mathbb EX\).
\end{hbtheorem}

Durrett states the first form as Exercise~1.5.5 and the second as
Exercise~1.6.13 \cite[pp.~28, 37]{DurrettPTEAuthor2019}; the source
writes \(g_n\uparrow g\) and \(X_n\uparrow X\) for the two monotone
hypotheses. In the example \(r_n\) above,
\(\int r_1^-\,d\lambda=\lambda([1,\infty))=\infty\), so it is
excluded by the hypothesis, and the conclusion fails for it.

\begin{hbtheorem}{thm:p15-w02c-p15-005-signed-series}{proved=annex:p15-w02c-p15-005-signed-series}{Term-by-term integration of a signed series}
\par\nobreak
Let \(f_1,f_2,\ldots\) be real measurable functions with
\(\sum_n\int|f_n|\,d\mu<\infty\). Then
\(G\defeq\sum_n|f_n|\) is integrable, so \(G<\infty\) a.e., and
\(\sum_nf_n(\omega)\) converges absolutely at every \(\omega\) with
\(G(\omega)<\infty\). Let \(F\) be that sum where \(G<\infty\) and
zero where \(G=\infty\). Then \(F\) is measurable and integrable,
\(\sum_n\left|\int f_n\,d\mu\right|<\infty\), and
\[
 \sum_n\int f_n\,d\mu=\int\sum_nf_n\,d\mu\defeq\int F\,d\mu .
\]
For extended real integrable \(f_n\), the same holds after each
\(f_n\) is replaced by the real-valued representative of
\cref{def:p15-integrable-functions}.
\end{hbtheorem}

This is Durrett's Exercise~1.5.10 \cite[p.~28]{DurrettPTEAuthor2019}.
The case \(f_n=f\mathbf1_{A_n}\) with disjoint \(A_n\) is the countable
additivity of \cref{thm:p15-signed-integral-laws}, and the nonnegative
case without any summability condition is part of
\cref{thm:p15-nonnegative-integral-laws}.

\begin{hbtheorem}{thm:p15-w02c-p15-005-indefinite-continuity}{proved=annex:p15-w02c-p15-005-indefinite-continuity}{Continuity of an indefinite integral}
\par\nobreak
Let \(a<b\) be real and let \(f\) be a real Borel function with
\(\int_{[a,b]}|f|\,d\lambda<\infty\). Then
\[
 g(x)\defeq\int_{[a,x]}f(y)\,dy\defeq\int f\mathbf1_{[a,x]}\,d\lambda,
 \qquad a\leq x\leq b,
\]
is continuous on \([a,b]\), in particular on \((a,b)\).
\end{hbtheorem}

Durrett's Exercise~1.5.8 asks for continuity on \((a,b)\)
\cite[p.~28]{DurrettPTEAuthor2019}; the proof also gives one-sided
continuity at the two endpoints.

\section{Almost-everywhere convergence and convergence in measure}
\label{sec:p15-w02c-p15-005-modes}

\begin{hbtheorem}{thm:p15-w02c-p15-005-ae-measure}{proved=annex:p15-w02c-p15-005-ae-measure}{Almost-everywhere convergence implies convergence in measure on a finite measure space}
\par\nobreak
Let \(f_n,f\) be real measurable functions, and let convergence in
measure be that of \cref{def:p15-lp-convergence}.
\begin{enumerate}
\item If \(\mu(\Omega)<\infty\) and \(f_n\to f\) a.e., then
      \(f_n\to f\) in measure.
\item Without \(\mu(\Omega)<\infty\) the implication fails: on
      \((\mathbb R,\lambda)\), \(f_n=\mathbf1_{(n,n+1]}\) converges
      to zero at every point, while
      \(\lambda(\{|f_n|>1/2\})=1\) for every \(n\).
\item The converse fails even on a probability space: the typewriter
      sequence \(T_n=\mathbf1_{[j2^{-k},(j+1)2^{-k})\cap(0,1)}\),
      \(n=2^k+j\) with \(0\leq j<2^k\), on the uniform probability
      space converges to zero in probability and at no point.
\end{enumerate}
\end{hbtheorem}

Durrett remarks on page~26 that on a space of finite measure
convergence in measure is a weaker assumption than a.e.\ convergence
\cite[p.~26]{DurrettPTEAuthor2019}. Part~(1) is the implication behind
that remark, on every finite measure space, and part~(3) shows that the
weakening is strict; the typewriter sequence is
\cref{thm:p15-w02b-p15-013-countermodels}~(2), where both of its
properties are proved. On a probability space part~(1) is
\cref{thm:p15-w02b-p15-013-as-probability}~(1).

Durrett makes the remark as he introduces his bounded convergence
theorem, Theorem~1.5.3, whose hypothesis is convergence in measure
\cite[p.~26]{DurrettPTEAuthor2019}; here that theorem is
\cref{thm:p15-dominated-bounded-convergence}~(2). Part~(1) shows that
an a.e.\ hypothesis is a special case of it.

Durrett gives the same name to a second statement. His Theorem~1.6.7,
the dominated convergence theorem for expectations, assumes
\(X_n\to X\) a.s., \(|X_n|\leq Y\) for every \(n\) and
\(\mathbb EY<\infty\), and concludes \(\mathbb EX_n\to\mathbb EX\); he
calls its special case with a constant \(Y\) the bounded convergence
theorem \cite[thm.~1.6.7, p.~30]{DurrettPTEAuthor2019}. Theorem~1.6.7
is \cref{thm:p15-dominated-bounded-convergence}~(1) on a probability
space. Its constant case is the corollary below with \(\mu=\mathbb P\)
and \(E=\Omega\). By part~(1) it is also a special case of
Theorem~1.5.3. The two statements that Durrett calls bounded
convergence therefore differ only in the mode of convergence they
assume, and on a finite measure space the version with convergence in
measure contains the almost-everywhere version.

\begin{hbcorollary}{thm:p15-w02c-p15-005-ae-bounded}{proved=annex:p15-w02c-p15-005-ae-bounded}{Bounded convergence from almost-everywhere convergence}
\par\nobreak
Let \(E\in\mathcal F\) with \(\mu(E)<\infty\), let \(0\leq M<\infty\),
and let \(f_n,f\) be real measurable functions such that each \(f_n\)
vanishes outside \(E\) a.e., \(|f_n|\leq M\) a.e., and \(f_n\to f\)
a.e. Then \(f_n\to f\) in measure, \(\|f_n-f\|_1\to0\), and
\(\int f_n\,d\mu\to\int f\,d\mu\). The total measure
\(\mu(\Omega)\) may be infinite.
\end{hbcorollary}

The condition \(\mu(E)<\infty\) cannot be dropped even for convergence
that is uniform on the whole line. Durrett's Example~1.5.4 takes
\(f_n=1/n\) on \([0,n]\) and zero elsewhere, on the real line with
Lebesgue measure \cite[p.~26]{DurrettPTEAuthor2019}; it differs only
at the point zero from the sequence \(h_n\) of
\cref{def:p15-convergence-failures}, which is bounded by one, converges
uniformly to zero, and has every integral equal to one.

\section{Jensen's inequality in several forms}
\label{sec:p15-w02c-p15-005-jensen}

\begin{hbexample}{def:p15-w02c-p15-005-jensen-direction}{Which way Jensen's inequality points, and Durrett's Figure~1.6}
Let \(0<\lambda<1\), let \(x,y\) be real, and let
\(\mathbb P(X=x)=\lambda\) and \(\mathbb P(X=y)=1-\lambda\). For
convex \(\varphi\),
\[
 \mathbb E\varphi(X)=\lambda\varphi(x)+(1-\lambda)\varphi(y)
 \geq\varphi(\lambda x+(1-\lambda)y)=\varphi(\mathbb EX).
\]
For a two-point law, Jensen's inequality is the defining inequality of
convexity, read with the weights as probabilities: the average of the
values lies above the value at the average.

Durrett's Figure~1.6 draws the case \(g(x)=x^2-3x+3\) with
\(\mathbb P(X=1)=\mathbb P(X=3)=1/2\). This \(g\) is convex: its
affine part satisfies the convexity relation with equality, and
\(tx^2+(1-t)y^2-(tx+(1-t)y)^2=t(1-t)(x-y)^2\geq0\). Here
\(\mathbb EX=2\), \(g(1)=1\), \(g(3)=3\) and \(g(2)=1\), so
\[
 \mathbb Eg(X)=\tfrac12\cdot1+\tfrac12\cdot3=2
 \;>\;1=g(\mathbb EX).
\]
In the figure, \(\mathbb Eg(X)\) is the midpoint of the chord from
\((1,1)\) to \((3,3)\), which lies above the graph at \(x=2\). The
inequality is strict, as \cref{thm:p15-w02c-p15-005-jensen-equality}
predicts, because \(g\) is strictly convex and \(X\) is not a.s.\ constant.
\end{hbexample}

The illustration and the figure are on Durrett's pages~29
and~30 \cite[pp.~29--30]{DurrettPTEAuthor2019}.

\begin{hbdefinition}{def:p15-w02c-p15-005-convex-rn}{Convex and strictly convex functions on \(\mathbb R^n\)}
A function \(\varphi:\mathbb R^n\to\mathbb R\) is \emph{convex} when
\(\varphi(tx+(1-t)y)\leq t\varphi(x)+(1-t)\varphi(y)\) for all
\(x,y\in\mathbb R^n\) and \(0<t<1\). It is \emph{strictly convex}
when the inequality is strict whenever \(x\neq y\). For \(n=1\) the
first notion is \cref{def:p15-convex-integral}. For
\(m,v\in\mathbb R^n\), \(m\cdot v\) is the coordinate sum
\(\sum_im_iv_i\).
\index{convex function!on Euclidean space}\index{strictly convex function}
\end{hbdefinition}

For \(n=1\), strict convexity is the hypothesis of Durrett's
Exercise~1.6.1, which asks that the convexity inequality hold strictly
for \(0<t<1\), and convexity on \(\mathbb R^n\) is the hypothesis of
his Exercise~1.6.2 \cite[p.~35]{DurrettPTEAuthor2019}. The lemma below
supplies the continuity and the supporting affine function that those
exercises use; for \(n=1\) it repeats, with the strict form added, what
the proof of \cref{thm:p15-jensen-integral} establishes for
\cref{def:p15-convex-integral}.

\begin{hblemma}{thm:p15-w02c-p15-005-supporting-affine}{proved=annex:p15-w02c-p15-005-supporting-affine}{Continuity and a supporting affine function on \(\mathbb R^n\)}
\par\nobreak
Let \(\varphi:\mathbb R^n\to\mathbb R\) be convex. Then \(\varphi\) is
continuous, hence Borel measurable. For every \(c\in\mathbb R^n\) there
is \(m\in\mathbb R^n\) with
\[
 \varphi(x)\geq\varphi(c)+m\cdot(x-c)\qquad\text{for all }x\in\mathbb R^n .
\]
If \(\varphi\) is strictly convex, this inequality is strict for every
\(x\neq c\).
\end{hblemma}

\begin{hbtheorem}{thm:p15-w02c-p15-005-jensen-several}{proved=annex:p15-w02c-p15-005-jensen-several}{Jensen's inequality for several variables}
\par\nobreak
Let \(\varphi:\mathbb R^n\to\mathbb R\) be convex, and let
\(X_1,\ldots,X_n\) be real random variables with
\(\mathbb E|X_i|<\infty\) for every \(i\). Then
\(\varphi(X_1,\ldots,X_n)\) is a random variable, and if
\(\mathbb E|\varphi(X_1,\ldots,X_n)|<\infty\),
\[
 \mathbb E\varphi(X_1,\ldots,X_n)\geq
 \varphi(\mathbb EX_1,\ldots,\mathbb EX_n).
\]
\end{hbtheorem}

This is Durrett's Exercise~1.6.2, which asks for an imitation of the
supporting-line proof of his Theorem~1.5.1
\cite[pp.~25, 35]{DurrettPTEAuthor2019}. For \(n=1\) it is
\cref{thm:p15-jensen-integral}, Durrett's Theorem~1.6.2
\cite[p.~29]{DurrettPTEAuthor2019}.

\begin{hbcorollary}{thm:p15-w02c-p15-005-jensen-special}{proved=annex:p15-w02c-p15-005-jensen-special}{Two special cases of Jensen's inequality}
\par\nobreak
For a real random variable \(X\) with \(\mathbb E|X|<\infty\),
\[
 |\mathbb EX|\leq\mathbb E|X|,
 \qquad
 (\mathbb EX)^2\leq\mathbb E(X^2),
\]
where the right side of the second inequality may be \(\infty\).
\end{hbcorollary}

Durrett names these two cases directly after his Theorem~1.6.2
\cite[p.~29]{DurrettPTEAuthor2019}. The first also follows from the
absolute-value bound of \cref{thm:p15-signed-integral-laws}, and the
second from the variance identity of \cref{thm:p15-expectation-bounds};
the proof below derives both from Jensen's inequality, as the source
presents them.

\begin{hbtheorem}{thm:p15-w02c-p15-005-jensen-equality}{proved=annex:p15-w02c-p15-005-jensen-equality}{The equality case of Jensen's inequality}
\par\nobreak
Let \(\varphi:\mathbb R\to\mathbb R\) be strictly convex, and let
\(X\) be a real random variable with \(\mathbb E|X|<\infty\) and
\(\mathbb E|\varphi(X)|<\infty\). Then
\(\varphi(\mathbb EX)=\mathbb E\varphi(X)\) if and only if
\(X=\mathbb EX\) a.s.
\end{hbtheorem}

The forward implication is Durrett's Exercise~1.6.1
\cite[p.~35]{DurrettPTEAuthor2019}; the reverse implication is
immediate and is included so that the statement characterizes equality.

\begin{hbtheorem}{thm:p15-w02c-p15-005-am-gm}{proved=annex:p15-w02c-p15-005-am-gm}{The weighted arithmetic--geometric mean inequality}
\par\nobreak
Let \(p(1),\ldots,p(n)>0\) with \(\sum_{m=1}^np(m)=1\), and let
\(y_1,\ldots,y_n>0\). Then
\[
 \sum_{m=1}^np(m)y_m\geq\prod_{m=1}^ny_m^{p(m)},
\]
with equality exactly when \(y_1=\cdots=y_n\). When \(p(m)=1/n\),
the arithmetic mean of \(y_1,\ldots,y_n\) is at least their geometric
mean.
\end{hbtheorem}

The inequality is Durrett's Exercise~1.6.12, obtained from Jensen's
inequality with \(\varphi(x)=e^x\) and
\(\mathbb P(X=\log y_m)=p(m)\) \cite[p.~37]{DurrettPTEAuthor2019};
the equality clause applies \cref{thm:p15-w02c-p15-005-jensen-equality}.

\section{What moments say about tail probabilities}
\label{sec:p15-w02c-p15-005-moment-bounds}

The squared bound \(a^2\,\mathbb P(|X|\geq a)\leq\mathbb EX^2\), which
is Durrett's display~(1.6.1), and the two usages of the names Chebyshev
and Markov that he reports for it are set out in
\cref{def:p15-w02c-p15-001-chebyshev-names,thm:p15-w02c-p15-001-chebyshev-square}.
The results below use that bound for \(a>0\), where it is the case
\(\varphi(x)=x^2\), \(A=\{x:|x|\geq a\}\) of
\cref{thm:p15-expectation-bounds}.

\begin{hbcorollary}{thm:p15-w02c-p15-005-chebyshev-displays}{proved=annex:p15-w02c-p15-005-chebyshev-displays}{The variance is at most the second moment}
\par\nobreak
For a real random variable \(X\) with \(\mathbb EX^2<\infty\),
\(\operatorname{Var}(X)\leq\mathbb EX^2\), with equality exactly when
\(\mathbb EX=0\).
\end{hbcorollary}

This is Durrett's display~(1.6.3), which he draws from the variance
formula~(1.6.2) \cite[p.~33]{DurrettPTEAuthor2019}.

\begin{hbtheorem}{thm:p15-w02c-p15-005-chebyshev-sharp}{proved=annex:p15-w02c-p15-005-chebyshev-sharp}{Chebyshev's bound is attained, and is not sharp in the limit}
\par\nobreak
\begin{enumerate}
\item If \(0<b\leq a\), the law
      \(\mathbb P(X=a)=b^2/a^2\), \(\mathbb P(X=0)=1-b^2/a^2\) has
      \(\mathbb EX^2=b^2\) and \(\mathbb P(|X|\geq a)=b^2/a^2\). So
      equality holds in \(a^2\mathbb P(|X|\geq a)\leq\mathbb EX^2\).
\item If \(0<\mathbb EX^2<\infty\), then
      \(\lim_{a\to\infty}a^2\,\mathbb P(|X|\geq a)/\mathbb EX^2=0\).
\end{enumerate}
\end{hbtheorem}

This is Durrett's Exercise~1.6.3 \cite[p.~35]{DurrettPTEAuthor2019}.
So no constant smaller than one can replace the factor in the bound for
all laws at a fixed \(a\), but for each fixed law the bound is eventually
far from equality.

\begin{hbtheorem}{thm:p15-w02c-p15-005-one-sided}{proved=annex:p15-w02c-p15-005-one-sided}{The one-sided Chebyshev bound}
\par\nobreak
\begin{enumerate}
\item Let \(a>0\), \(b>0\) and \(0<p<1\), and let \(X\) have
      \(\mathbb P(X=a)=p\) and \(\mathbb P(X=-b)=1-p\). If \(Y\) is a
      random variable with \(\mathbb EY^2<\infty\),
      \(\mathbb EY=\mathbb EX\) and
      \(\operatorname{Var}(Y)=\operatorname{Var}(X)\), then
      \(\mathbb P(Y\geq a)\leq p\), with equality when \(Y=X\).
\item If \(\mathbb EY=0\), \(\operatorname{Var}(Y)=\sigma^2<\infty\)
      and \(a>0\), then
      \[
       \mathbb P(Y\geq a)\leq\frac{\sigma^2}{a^2+\sigma^2},
      \]
      and for each such \(a\) and \(\sigma\geq0\) some \(Y\) attains
      equality.
\end{enumerate}
\end{hbtheorem}

This is Durrett's Exercise~1.6.4 \cite[p.~36]{DurrettPTEAuthor2019}.
The source states part~(1) for \(a>b>0\), which is included; the proof
uses only \(a>0\) and \(b>0\), and part~(2) needs the extra range
\(b\geq a\) when \(a\leq\sigma\). Part~(2) improves the two-sided bound
\(\sigma^2/a^2\) for the upper tail alone and, unlike it, never
exceeds one.

\begin{hbtheorem}{thm:p15-w02c-p15-005-no-lower}{proved=annex:p15-w02c-p15-005-no-lower}{A mean and a variance give no lower bound on a tail}
\par\nobreak
\begin{enumerate}
\item For every \(\varepsilon>0\),
      \(\inf\{\mathbb P(|X|>\varepsilon):\mathbb EX=0,\
      \operatorname{Var}(X)=1\}=0\).
\item For every \(y\geq1\) and \(\sigma^2\in(0,\infty)\),
      \(\inf\{\mathbb P(|X|>y):\mathbb EX=1,\
      \operatorname{Var}(X)=\sigma^2\}=0\).
\end{enumerate}
The infima range over real random variables with finite second moment
and the stated mean and variance.
\end{hbtheorem}

This is Durrett's Exercise~1.6.5 \cite[p.~36]{DurrettPTEAuthor2019}.

\begin{hbtheorem}{thm:p15-w02c-p15-005-positive-lower}{proved=annex:p15-w02c-p15-005-positive-lower}{A lower bound for the probability of a positive value}
\par\nobreak
Let \(Y\geq0\) with \(\mathbb EY^2<\infty\). Then
\((\mathbb EY)^2\leq\mathbb EY^2\cdot\mathbb P(Y>0)\). Hence, if
\(\mathbb EY^2>0\),
\[
 \mathbb P(Y>0)\geq\frac{(\mathbb EY)^2}{\mathbb EY^2} .
\]
If \(\mathbb EY^2=0\), then \(Y=0\) a.s., \(\mathbb P(Y>0)=0\), and
the displayed quotient has no value.
\end{hbtheorem}

This is Durrett's Exercise~1.6.6, proved, as the source directs, by
the Cauchy--Schwarz inequality applied to \(Y\mathbf1_{(Y>0)}\)
\cite[p.~36]{DurrettPTEAuthor2019}. That inequality is the case
\(p=q=2\) of \cref{thm:p15-holder-integral}; Durrett's direct proof of
it from a nonnegative quadratic is
\cref{thm:p15-w02c-p15-001-cauchy-schwarz}.

\begin{hbtheorem}{thm:p15-w02c-p15-005-no-dominator}{proved=annex:p15-w02c-p15-005-no-dominator}{Moment control without an integrable dominator}
\par\nobreak
On the uniform probability space \(((0,1),\mathcal F_U,\mathbb U)\),
let \(1<\alpha<2\) and
\(X_n=n^\alpha\mathbf1_{(1/(n+1),1/n)}\). Then:
\begin{enumerate}
\item \(X_n\to0\) at every point, hence a.s.;
\item with \(h(x)=x\) and \(g(x)=|x|^{2/\alpha}\), the hypotheses of
      \cref{thm:p15-moment-controlled-convergence} hold with \(K=1\),
      so \(\mathbb EX_n\to0\); directly,
      \(\mathbb EX_n=n^{\alpha-1}/(n+1)\);
\item no integrable \(Y\) satisfies \(|X_n|\leq Y\) a.s.\ for every
      \(n\).
\end{enumerate}
\end{hbtheorem}

This is Durrett's Exercise~1.6.7 \cite[p.~36]{DurrettPTEAuthor2019};
the theorem it refers to is his Theorem~1.6.8
\cite[p.~31]{DurrettPTEAuthor2019}, which is
\cref{thm:p15-moment-controlled-convergence}. So the moment-control
theorem reaches limits that dominated convergence
(\cref{thm:p15-dominated-bounded-convergence}) cannot.

\begin{hbtheorem}{thm:p15-w02c-p15-005-inverse-tail}{proved=annex:p15-w02c-p15-005-inverse-tail}{The inverse moment beyond a level}
\par\nobreak
Let \(X\geq0\) be a real random variable; \(\mathbb E(1/X)<\infty\)
is not assumed. For \(y>0\), \(\mathbb E(1/X;X>y)\leq1/y\), and
\[
 \lim_{y\to\infty}y\,\mathbb E(1/X;X>y)=0,
 \qquad
 \lim_{y\downarrow0}y\,\mathbb E(1/X;X>y)=0 .
\]
\end{hbtheorem}

This is Durrett's Exercise~1.6.14 \cite[p.~37]{DurrettPTEAuthor2019}.
On \(\{X>y\}\) the integrand \(1/X\) is finite, so the restricted
expectation is that of a nonnegative random variable.

\section{Inclusion--exclusion and Bonferroni bounds}
\label{sec:p15-w02c-p15-005-bonferroni}

\begin{hbdefinition}{def:p15-w02c-p15-005-intersection-sums}{Intersection sums}
Let \(A_1,\ldots,A_n\) be events and \(A=\bigcup_{i=1}^nA_i\). For a
nonempty \(S\subseteq\{1,\ldots,n\}\), let
\(A_S\defeq\bigcap_{i\in S}A_i\). For \(1\leq k\leq n\), the
\(k\)-th \emph{intersection sum} is
\[
 S_k\defeq\sum_{|S|=k}\mathbb P(A_S)
 =\sum_{i_1<\cdots<i_k}\mathbb P(A_{i_1}\cap\cdots\cap A_{i_k}).
\]
\index{inclusion--exclusion}\index{Bonferroni inequalities}
\end{hbdefinition}

\begin{hbtheorem}{thm:p15-w02c-p15-005-inclusion-exclusion}{proved=annex:p15-w02c-p15-005-inclusion-exclusion}{The inclusion--exclusion formula}
\par\nobreak
With the notation of \cref{def:p15-w02c-p15-005-intersection-sums},
\(\mathbf1_A=1-\prod_{i=1}^n(1-\mathbf1_{A_i})\), and
\[
 \begin{aligned}
 \mathbb P\Bigl(\bigcup_{i=1}^nA_i\Bigr)
 &=\sum_{k=1}^n(-1)^{k-1}S_k\\
 &=\sum_i\mathbb P(A_i)-\sum_{i<j}\mathbb P(A_i\cap A_j)\\
 &\qquad+\sum_{i<j<k}\mathbb P(A_i\cap A_j\cap A_k)-\cdots
  +(-1)^{n-1}\mathbb P\Bigl(\bigcap_{i=1}^nA_i\Bigr).
 \end{aligned}
\]
\end{hbtheorem}

This is Durrett's Exercise~1.6.9, proved by expanding the product and
taking expected values, as the source directs
\cite[p.~36]{DurrettPTEAuthor2019}.

\begin{hbtheorem}{thm:p15-w02c-p15-005-bonferroni}{proved=annex:p15-w02c-p15-005-bonferroni}{The Bonferroni inequalities}
\par\nobreak
With the notation of \cref{def:p15-w02c-p15-005-intersection-sums},
for each odd \(m\) with \(1\leq m\leq n\),
\[
 \mathbb P\Bigl(\bigcup_{i=1}^nA_i\Bigr)\leq\sum_{k=1}^m(-1)^{k-1}S_k ,
\]
and for each even \(m\) with \(2\leq m\leq n\),
\[
 \mathbb P\Bigl(\bigcup_{i=1}^nA_i\Bigr)\geq\sum_{k=1}^m(-1)^{k-1}S_k .
\]
The cases \(m=1,2,3\) are
\(\mathbb P(A)\leq S_1\), \(\mathbb P(A)\geq S_1-S_2\) and
\(\mathbb P(A)\leq S_1-S_2+S_3\). Pointwise,
\(\mathbf1_A\leq\sum_{i=1}^n\mathbf1_{A_i}\), and in general the
truncated alternating sum of indicators is at least \(\mathbf1_A\)
after an odd number of sums and at most \(\mathbf1_A\) after an even
number.
\end{hbtheorem}

This is Durrett's Exercise~1.6.10, whose three displayed cases and
closing parity rule are on his page~37
\cite[pp.~36--37]{DurrettPTEAuthor2019}: stopping after an even number
of sums gives a lower bound and after an odd number an upper bound. The
case \(m=n\) is the identity of
\cref{thm:p15-w02c-p15-005-inclusion-exclusion}.

\section{Procedure, worked example and limits}
\label{sec:p15-w02c-p15-005-procedure}

\begin{hbmethod}{def:p15-w02c-p15-005-procedure}{Choosing a limit theorem or a probability bound}
\begin{enumerate}
\item \emph{Passing a limit through an integral.} For an increasing
      sequence, check \(\int g_1^-\,d\mu<\infty\) and use
      \cref{thm:p15-w02c-p15-005-monotone-lower}. For a series, check
      \(\sum_n\int|f_n|\,d\mu<\infty\) and use
      \cref{thm:p15-w02c-p15-005-signed-series}. For a.s.\ convergence
      on a probability space with a uniform bound
      \(\mathbb Eg(X_n)\leq K\) on a faster-growing \(g\), use
      \cref{thm:p15-moment-controlled-convergence}; a dominator need not
      exist (\cref{thm:p15-w02c-p15-005-no-dominator}).
\item \emph{Changing the mode of convergence.} A.e.\ convergence gives
      convergence in measure only when \(\mu(\Omega)<\infty\)
      (\cref{thm:p15-w02c-p15-005-ae-measure}); for functions bounded
      by a constant and vanishing outside a set of finite measure it
      then gives convergence of the integrals
      (\cref{thm:p15-w02c-p15-005-ae-bounded}).
\item \emph{Upper bounds from moments.} For \(\mathbb P(|X|\geq a)\)
      use \(\mathbb EX^2/a^2\); for the one-sided tail of a centred
      variable use \(\sigma^2/(a^2+\sigma^2)\). Side conditions: \(a>0\),
      a finite second moment, and for the one-sided bound mean zero.
\item \emph{Lower bounds.} A mean and a variance alone give none
      (\cref{thm:p15-w02c-p15-005-no-lower}). For \(Y\geq0\) with
      \(0<\mathbb EY^2<\infty\), use \((\mathbb EY)^2/\mathbb EY^2\).
\item \emph{Unions.} Compute \(S_1,S_2,\ldots\) and stop at the first
      \(m\) whose truncated alternating sums bracket the probability
      tightly enough; odd \(m\) gives an upper bound, even \(m\) a lower
      bound, and \(m=n\) the exact value.
\end{enumerate}
\end{hbmethod}

\begin{hbexample}{def:p15-w02c-p15-005-worked}{Moment bounds and Bonferroni bounds on small laws}
Let \(\mathbb P(Y=0)=\mathbb P(Y=2)=1/4\) and \(\mathbb P(Y=1)=1/2\).
Then \(\mathbb EY=1\), \(\mathbb EY^2=1/2+1=3/2\) and
\(\operatorname{Var}(Y)=1/2\).
\begin{itemize}
\item \(\mathbb P(Y>0)=3/4\), while
      \((\mathbb EY)^2/\mathbb EY^2=2/3\).
\item With \(a=2\): \(4\,\mathbb P(|Y|\geq2)=1\leq3/2\).
\item For \(Z=Y-1\), with mean zero and \(\sigma^2=1/2\), at \(a=1\):
      \(\mathbb P(Z\geq1)=1/4\leq(1/2)/(1+1/2)=1/3\), while the
      two-sided bound gives \(\mathbb P(|Z|\geq1)\leq1/2\), with
      equality, since \(\mathbb P(|Z|\geq1)=1/2\).
\end{itemize}
For the union bounds, roll a fair die and let
\(A_1=\{2,4,6\}\), \(A_2=\{4,5,6\}\) and \(A_3=\{1,6\}\). Then
\(S_1=1/2+1/2+1/3=4/3\); the pairwise intersections are
\(\{4,6\}\), \(\{6\}\), \(\{6\}\), so \(S_2=1/3+1/6+1/6=2/3\); and
\(S_3=\mathbb P(\{6\})=1/6\). The union is \(\{1,2,4,5,6\}\), with
probability \(5/6\). The bounds read \(5/6\leq4/3\),
\(5/6\geq4/3-2/3=2/3\), and
\(5/6=4/3-2/3+1/6\).
\end{hbexample}

\begin{hbwarning}{def:p15-w02c-p15-005-limits}{What fails without the hypotheses}
\begin{enumerate}
\item \emph{Without absolute summability} the series identity fails.
      On \((\mathbb R,\lambda)\) let
      \(f_n=\mathbf1_{(n-1,n]}-\mathbf1_{(n,n+1]}\). Each
      \(\int f_n\,d\lambda=0\) and \(\int|f_n|\,d\lambda=2\). The
      partial sums are
      \(\sum_{n\leq N}f_n=\mathbf1_{(0,1]}-\mathbf1_{(N,N+1]}\), which
      converge at every point to \(\mathbf1_{(0,1]}\), with integral
      one, while \(\sum_n\int f_n\,d\lambda=0\).
\item \emph{Without an integrable negative part} monotone convergence
      fails, by the example \(r_n\) recalled before
      \cref{thm:p15-w02c-p15-005-monotone-lower}.
\item \emph{Without strict convexity} the equality case of Jensen's
      inequality fails: for affine \(\varphi\), equality holds for
      every integrable \(X\).
\item \emph{Without the sign condition} the lower bound for
      \(\mathbb P(Y>0)\) fails: if \(Y=-1\) a.s., then
      \(\mathbb P(Y>0)=0\) while \((\mathbb EY)^2/\mathbb EY^2=1\).
\end{enumerate}
\end{hbwarning}

\emph{Relationships.} The monotone and series results extend
\cref{thm:p15-monotone-integral,thm:p15-nonnegative-integral-laws} to
signed functions; the a.e.-to-measure implication is used with
\cref{thm:p15-dominated-bounded-convergence}~(2), which assumes
convergence in measure; the several-variable Jensen inequality and its
equality case extend \cref{thm:p15-jensen-integral}, whose
conditional form is \cref{thm:p15-conditional-jensen}; and the
moment bounds refine \cref{thm:p15-expectation-bounds}.

\emph{Sources.} All from the author version dated 11 January 2019 of
Durrett's \emph{Probability: Theory and Examples}: Jensen's inequality
for integrals and H\"older's inequality, Theorems~1.5.1 and~1.5.2
\cite[p.~25]{DurrettPTEAuthor2019}; the remark on convergence in
measure, Theorem~1.5.3 and Example~1.5.4
\cite[p.~26]{DurrettPTEAuthor2019};
Exercises~1.5.5, 1.5.8 and~1.5.10 \cite[p.~28]{DurrettPTEAuthor2019};
Theorem~1.6.2 with its two-point illustration and special cases
\cite[p.~29]{DurrettPTEAuthor2019}; Figure~1.6
\cite[p.~30]{DurrettPTEAuthor2019}; Theorem~1.6.8
\cite[p.~31]{DurrettPTEAuthor2019}; displays~(1.6.2) and~(1.6.3)
\cite[p.~33]{DurrettPTEAuthor2019}; Exercises~1.6.1--1.6.3
\cite[p.~35]{DurrettPTEAuthor2019}; Exercises~1.6.4--1.6.7, 1.6.9
and the first part of~1.6.10 \cite[p.~36]{DurrettPTEAuthor2019}; and
the rest of Exercise~1.6.10 with Exercises~1.6.12--1.6.14
\cite[p.~37]{DurrettPTEAuthor2019}.
